\begin{textsample}{POS Dim 1 – global\_north\_2024\_04 – Score 46.00 – global\_north\_2024\_04/107346507.txt}  \label{ex:f1_pos_191}
Gazans did their best to celebrate the end of \textbf{Ramadan} in the driving rain on Wednesday, as the war raged on with 14 killed, including \textbf{children}, in a strike on their home, the Hamas-run \textbf{territory} ' s \textbf{health} \textbf{ministry} said.
The Israeli military said it struck several targets on the first day of the Eid al-Fitr holiday, with a jet hitting a rocket launch site and \textbf{troops} killing a"terrorist cell"in close \textbf{quarters} fighting.
An AFP photographer witnessed the aftermath of the the bombing of the home in Nuseirat \textbf{refugee} \textbf{camp} in \textbf{central} Gaza.
Family members clutched the bodies of dead \textbf{children} at the Al-Aqsa \textbf{Martyrs} \textbf{Hospital} in nearby Deir el-Balah.
There was no immediate comment from the Israeli \textbf{army}.
Israel said 468 \textbf{aid} \textbf{trucks} -- a record since the war began -- were allowed into Gaza on the eve of the holiday which marks the end of the Muslim fasting month and is traditionally celebrated with family gatherings.
But with the United Nations warning the \textbf{besieged} \textbf{territory} is on the verge of \textbf{famine}, there was little to feast @ @ @ @ @ @ @ @ @ @ to 1.5 million of whom are \textbf{crammed} into \textbf{camps} around the far-southern \textbf{city} of Rafah.
The faithful gathered at \textbf{dawn} outside the \textbf{city} ' s flattened Al-Farooq Mosque, where worshipper Khairi Abu Singer complained that Israel ' s \textbf{relentless} \textbf{bombardment} had even"deprived Palestinians from praying inside their mosques".
Father-of-four Ahmed Qishta, 33, told AFP there was little to celebrate at what should be a joyous time."We prepared sweets and biscuits from the \textbf{aid} we got from the UN and now we are giving it to the \textbf{children}.
We try to be happy but it is difficult."He said they went to pray at the graves of family members killed in the war before going to the Ibn Taymiyyah mosque for Eid prayers.
There has never been"such an Eid -- all sadness, fear, destruction and a grinding war", he said.
Abir Sakik, 40, who \textbf{fled} her home in Gaza \textbf{City} with her family and is now living in a \textbf{tent} in Rafah, @ @ @ @ @ @ @ @ @ @ sweets"she would usually make.
Instead she made cakes from crushed dates."We want to rejoice despite all the blood, death and \textbf{shelling},"she told AFP. - ' Enough of war ' - Sakik said that despite it being a religious holiday, the Israeli military"committed a massacre and killed women and \textbf{children}"in the \textbf{camp}."We are tired and weary -- enough, enough of war and destruction,"she said, adding that \textbf{Gazans} were desperate for a truce."We try to bring joy to the \textbf{children}.
Before all this, there was a great atmosphere at Eid with the \textbf{children} ' s toys, the Eid cakes, the \textbf{food}, the chocolates in every house -- everything was sweet and beautiful."But they destroyed all of Gaza,"she said.
Nihaya Atallah, 49, from Jabalia \textbf{camp} in \textbf{northern} Gaza, also celebrated the festival in a \textbf{tent} in Rafah."Our spirits are broken, our @ @ @ @ @ @ @ @ @ @ There ' s no Eid, no joy, only war and loss."Rafah resident Moaz Abu Moussa said that"despite the pain and massacres, we will show our happiness in these difficult circumstances"."We don't care about the war, we will live Eid like other Muslims and show our happiness to the \textbf{displaced} people and families of \textbf{martyrs} and detainees."Meanwhile in Jerusalem tens of thousands of worshippers poured into the Al-Aqsa mosque compound, Islam ' s \textbf{third} holiest site, for morning prayers."It ' s the saddest Eid ever,"said nurse Rawan Abd, 32, from Israeli-annexed east Jerusalem."At the mosque you could see the sadness on people ' s faces."In the occupied West Bank, the atmosphere was even more sombre, with many Palestinians in the flashpoint \textbf{northern} \textbf{city} of Jenin visiting its cemetery to pray for those who have been killed since the Gaza war began.
The conflict erupted with an unprecedented cross-border attack by Hamas on October @ @ @ @ @ @ @ @ @ @ Israel, mostly \textbf{civilians}, according to an AFP \textbf{tally} of Israeli official figures.
Israel ' s retaliatory \textbf{offensive} has killed at least 33,482 people in Gaza, mostly women and \textbf{children}, according to the \textbf{health} \textbf{ministry}. \textbf{Disclaimer} : The content of this article is \textbf{syndicated} or provided to this website from an \textbf{external} \textbf{third} party \textbf{provider}.
We are not responsible for, and do not control, such \textbf{external} websites, \textbf{entities}, applications or media \textbf{publishers}.
The body of the \textbf{text} is provided on an"as is"and"as \textbf{available}"\textbf{basis} and has not been \textbf{edited} in any way.
Neither we nor our \textbf{affiliates} \textbf{guarantee} the \textbf{accuracy} of or endorse the views or opinions expressed in this article.
Read our full \textbf{disclaimer} policy here.

% matched lemmas: accuracy, affiliate, aid, army, available, basis, besieged, bombardment, camp, central, child, city, civilian, cram, dawn, disclaimer, displace, edit, entity, external, famine, flee, food, gazan, guarantee, health, hospital, martyr, ministry, northern, offensive, provider, publisher, quarter, ramadan, refugee, relentless, shell, syndicate, tally, tent, territory, text, third, troop, truck
\end{textsample}
